
	
	\section{abstract}
		This document compares the established string theory with an alternative approach, the Fractal Field Geometry Theory (FFGFT). While both models reject the concept of point-like particles, they differ fundamentally in their treatment of dimensions, complexity and time. FFGFT proposes a four-dimensional, fractally self-similar torus field in which time is not a dimension but a local ratio to mass.

	
	\section*{Introduction: The Common Root}
	
	Both the established string theory and the model of Fractal Field Geometry Theory (FFGFT) presented here share a fundamental insight: the concept of the point-like particle as the fundamental building block of reality is inadequate. The actual foundation of the world must be an extended, geometric and self-contained object. But while string theory drives this idea into dimensions, FFGFT takes the path of fractal nesting.
	
	\section{The Foundation: String vs.\ Fractal Torus}
	
	\begin{itemize}
		\item \textbf{String theory:} Its foundation is a one-dimensional, vibrating object -- the string. This exists \textit{in} a space and its vibrational modes determine the properties of the particles arising from it. The complexity of the world arises through the multitude of possible vibrations in a high-dimensional space.
		
		\item \textbf{FFGFT (Torus model):} Its foundation is the field itself. There is no ``object in space''. The world is a single, continuous, four-dimensional field. A particle, an atom, a planet -- all of these are not separate objects, but merely local \textbf{torsion configurations} of this field: compressed vortices or nodes, comparable to a corkscrew pattern in the structure of space.
	\end{itemize}
	
	The \textbf{torus} (more precisely: a four-dimensional torus) is here the most stable and natural form for such a closed system. Its topology enables the field to loop back into itself -- without beginning, without end and without an ``outside''.
	
	\section{The Origin of Complexity: Dimensions vs.\ Self-Similarity}
	
	Here lies the most radical difference between the two approaches.
	
	\begin{table}[h]
		\centering
		\adjustbox{max width=\textwidth}{%
\begin{tabular}{p{4cm} p{5cm} p{5cm}}
			\toprule
			\textbf{Feature} & \textbf{String theory} & \textbf{FFGFT (Fractal Torus)} \\
			\midrule
			\textbf{Number of dimensions} & 10, 11 or more. The excess dimensions are rolled up into tiny, multi-dimensional forms (Calabi-Yau manifolds). & \textbf{Four dimensions}. No more are needed. \\
			\addlinespace
			\textbf{Source of complexity} & The infinite variety of vibration patterns and the geometric forms into which the extra dimensions can be rolled up. & \textbf{Fractal self-similarity}. The torus structure repeats itself at every scale: an electron vortex is a mini-torus, an atom a larger one, a galaxy an even larger one -- all follow the same geometric blueprint. \\
			\addlinespace
			\textbf{Predictability} & A ``landscape problem'': the theory permits approximately \textbf{10\textsuperscript{500}} different universes each with their own natural laws. It thereby loses its predictive power. & A single fundamental parameter: \textbf{$\xi = 4/30000$}. From it the fine structure constant $\alpha$ (the strength of the electromagnetic interaction) can be derived. From this in turn the masses of all particles and finally the gravitational constant $G$ follow. One parameter determines everything. \\
			\bottomrule
		\end{tabular}%
}
		\caption{Comparison between string theory and FFGFT}
		\label{tab:vergleich}
	\end{table}
	
	Complexity in FFGFT is not a matter of \textit{more} (dimensions), but of \textit{deeper} (nesting). A vortex generates smaller vortices, which in turn generate smaller ones -- an endless cascade of self-similarity reaching from the entire universe down to the smallest scales.
	
	\section{The Redefinition of Time: From Container to Local Ratio}
	
	In classical physics (and string theory), time is the fourth dimension -- a ``container'', a coordinate like height, width and depth. FFGFT breaks radically with this notion.
	
	\begin{itemize}
		\item \textbf{The fourth dimension of the torus is not time.} It is a \textit{topological} property -- the closure of the field. It is the condition for there being such a thing as an ``inside'' and ``outside'' at all, but is not itself a place one can travel to.
		
		\item \textbf{Time is a local ratio.} It is not a universal coordinate but a local property of the torsion density of the field:
		
		\begin{equation}
			\tilde{T} = \frac{1}{m}
			\label{eq:time}
		\end{equation}
		
		\textit{Where there is much mass (and thus torsion), time passes locally more slowly. Where there is no mass, there is no local time.} A photon is massless, so it ``experiences'' no time -- not because it moves at the speed of light, but because there is simply no ``inner clock'' (no torsion compression) within it.
		
		\item \textbf{Why we experience time linearly.} We are not external observers of time, we are part of the field. We are torsion configurations that vibrate along with the fractal unfolding of the torus. This unfolding is inherently directed (a fractal can only grow, not shrink). This directed change of our own structure \textit{is} what we experience as the flow of time. Time does not flow; \textit{we flow}.
	\end{itemize}
	
	\section{Finite but Endless: The Perspective from Within}
	
	The deepest and perhaps most consequential implication of the model is our position within it.
	
	\begin{quote}
		\textit{``We are not observers who contemplate the torus. We are the torus contemplating itself.''}
	\end{quote}
	
	The universe is \textbf{finite} -- it has a determinable volume. But it is \textbf{endless} -- it has no boundary, no wall one could reach. Like the surface of a sphere it is closed within itself; one can move on it forever without ever reaching an end. The difference is only that this closure takes place in four dimensions.
	
	For us as part of this structure this has a bold consequence: the question ``What came before the Big Bang?'' is meaningless. It presupposes a ``before'' and an ``outside'' that do not exist in a closed geometry. Just as the question ``What lies north of the North Pole?'' has no answer, because the geometry of the globe dissolves the question.
	
	\section*{Conclusion: String Theory and FFGFT Compared}
	
	String theory describes one-dimensional loops vibrating in a high-dimensional, complex space. Its approach is additive: more dimensions, more forms, more possibilities.
	
	FFGFT describes a four-dimensional field with fractal torsion structure. Its approach is integrative: from a single geometry and a single parameter the entire complexity of the world unfolds -- from quantum physics to cosmology.
	
	The core of the difference lies in the treatment of space and time:
	\begin{itemize}
		\item In string theory, time is a dimension. We move \textit{in} it.
		\item In FFGFT, time is a ratio. It is \textit{we} who change.
	\end{itemize}
	
	\vspace{1cm}
	\begin{center}
		\rule{5cm}{0.4pt}
		
		\small{Documented on \today}
	\end{center}
